% Targeted replacements for IEEE_WCL_SSFL (13).pdf
%
% Purpose:
% - Replace only the parts that still read awkwardly in version (13).
% - Leave untouched sections alone.
%
% Replace the following:
% 1) Fig. 1 caption and label
% 2) Entire Section II
% 3) Section III-A only
%
% Main differences from version (13):
% - Section II wording is smoothed, but the structure stays the same.
% - Section III-A is reduced back to a true overview subsection.
% - The long proposed data-flow equations in III-A are removed.
% - Only one summary equation is kept in III-A:
%     M_final = M_vib \odot M_ch
% - Detailed equations remain in the later subsections, where they belong.

% ------------------------------------------------------------------
% Fig. 1 caption replacement
% ------------------------------------------------------------------
% Replace the current Fig. 1 caption and remove any leftover Korean note.
%
% \caption{Generic split federated learning framework considered in this work. Multiple clients collaboratively train a split model by transmitting smashed data to the server over wireless uplink channels and receiving intermediate gradients for backpropagation.}
% \label{fig:system_model}

% ------------------------------------------------------------------
% Section II replacement
% ------------------------------------------------------------------
\section{System Model}

As illustrated in Fig.~\ref{fig:system_model}, we consider an SFL framework \cite{gupta2018distributed,thapa2022splitfed}. In the considered setting, clients communicate with the server over wireless uplink channels. The global model is divided into a client-side sub-model $\phi_c$ and a server-side sub-model $\phi_s$.

For an input sample $x$, the client first computes an intermediate feature representation, commonly referred to as smashed data, as
\begin{equation}
    f = \phi_c(x).
\end{equation}
The smashed data $f$ is then transmitted to the server through the wireless uplink. The received representation at the server is modeled as
\begin{equation}
    f' = h \cdot f + n,
\end{equation}
where $n \sim \mathcal{CN}(0,\sigma_n^2 I)$ denotes additive white Gaussian noise (AWGN) and $h$ denotes the channel coefficient. In particular, $h=1$ for the AWGN channel and $h \sim \mathcal{CN}(0,1)$ for the Rayleigh fading channel. Based on the received representation, the server performs the remaining layers of the split model to produce the final prediction
\begin{equation}
    \hat{y} = \phi_s(f').
\end{equation}

This split architecture reduces the computational burden on client devices by offloading most of the deep neural network computation to the server, which makes it attractive for resource-limited IoT environments. In the generic SFL setting, the split model is trained end-to-end using a standard supervised objective defined on the server output. For a ground-truth label $y$, the task loss is written as
\begin{equation}
    \mathcal{L}_{\mathrm{task}} = \ell(\hat{y}, y),
\end{equation}
where $\ell(\cdot,\cdot)$ denotes a classification loss such as cross-entropy. During training, the server computes the loss from the final prediction and propagates the intermediate gradient back through the split point so that both $\phi_s$ and $\phi_c$ are updated. After local updates, the client-side model parameters are aggregated to form the global model for the next communication round.

Although SFL is effective in reducing on-device computation, its communication cost remains substantial because high-dimensional smashed data must be transmitted at every training iteration \cite{gao2020end}. This problem becomes more severe when the client-server exchange takes place over wireless channels, where limited bandwidth and channel noise can significantly degrade transmission efficiency and downstream task performance. These challenges motivate the proposed channel-aware semantic transmission mechanism described in the next section.

% ------------------------------------------------------------------
% Section III-A replacement only
% ------------------------------------------------------------------
\subsection{Proposed CA-SSFL Architecture}
Building on the SFL framework in Section II, we propose a channel-aware semantic split federated learning (CA-SSFL) architecture to reduce communication overhead while preserving task-relevant information under dynamic channel conditions. Motivated by task-oriented semantic communication \cite{qin2021semantic,gunduz2022beyond}, the proposed method augments the conventional client-server split model with a semantic encoder at the client side and a semantic decoder at the server side, as illustrated in Fig.~\ref{fig:proposed_encoder}. Accordingly, the client-side model is defined as $\Theta_C = \{\phi_c, \theta_{en}\}$, where $\phi_c$ extracts smashed data and $\theta_{en}$ compresses it into a transmission-efficient latent representation. The server-side model is defined as $\Theta_S = \{\theta_{de}, \phi_s\}$, where $\theta_{de}$ reconstructs the received latent representation and $\phi_s$ performs the downstream prediction task. Therefore, the overall model is given by
\begin{equation}
\Theta = \{\phi_c, \theta_{en}, \theta_{de}, \phi_s\}.
\end{equation}

Unlike conventional SFL, the proposed CA-SSFL incorporates two key mechanisms in the client-side semantic encoder: 1) a VIB-based semantic masking module that estimates the task relevance of latent features, and 2) an SNR-driven channel masking module that determines their transmission feasibility under the current wireless condition. The final transmission support is determined by combining these two masks as
\begin{equation}
M_{\mathrm{final}} = M_{\mathrm{vib}} \odot M_{\mathrm{ch}},
\end{equation}
so that only the feature dimensions that are both semantically important and channel-feasible are selected for transmission. The selected latent features are then sent through a dynamic slicing operation, and the server restores the sparse representation to its original dimensionality before processing it through the semantic decoder and the remaining server-side network. The detailed operations of the semantic encoder, channel masking, dynamic slicing, and end-to-end training objective are presented in the following subsections.

% ------------------------------------------------------------------
% Still need manual cleanup in v13
% ------------------------------------------------------------------
% 1) Delete the visible placeholder text:
%    "(무선 SFL 그림으로 수정)"
%
% 2) Change the current Fig. 2 label/cross-reference if needed so that
%    Section III-A can use:
%    \label{fig:proposed_encoder}
%
% 3) Keep the already-corrected phrase in III-B:
%    "dimension-wise KL divergence"
%    Do not revert it to "channel-wise KL divergence".
